\documentclass[11pt]{article}
\input{../../shared/project_style.tex}
\rhead{Basics 32 Textbook Draft}

\begin{document}
\DocTitle{Basics 32: Alfvén Waves from First Principles}{V - Plasma waves and Alfvén physics}{A full linear ideal-MHD derivation of the shear-Alfvén wave and its physical energy balance}

\begin{overviewbox}
\textbf{Textbook draft, version 1.0.} The shear-Alfvén wave is the central elementary wave of magnetized fluid plasma. Starting from a uniform equilibrium, this chapter derives its dispersion relation, polarization, wave equation, energy partition, and propagation along magnetic field lines. This chapter is designed for a reader with no prior plasma-physics training. It distinguishes exact identities from physical approximations and connects the central equations to later simulation modules.
\end{overviewbox}

\ProjectTOC

\section{Learning objectives}
After completing this note, the reader should be able to:
\begin{itemize}[leftmargin=1.6em]
\item Identify the equilibrium and perturbation assumptions of a shear-Alfvén wave.
\item Derive linearized induction and momentum relations.
\item Derive $\omega^2=k_\parallel^2v_A^2$.
\item Interpret magnetic tension as the restoring force.
\item Relate velocity and magnetic perturbation amplitudes.
\item Derive kinetic-magnetic energy equipartition for a traveling wave.
\end{itemize}

\section{Prerequisites and reading strategy}
\begin{itemize}[leftmargin=1.6em]
\item Basics 24, 25, 29, and 31.
\end{itemize}
On a first reading, emphasize physical meaning and dimensional checks. On a second reading, reproduce the derivations. Attempt the computational laboratory only after the limiting cases are understood.

\section{Uniform equilibrium and shear polarization}

Take
\begin{equation}
\mathbf B_0=B_0\hat{\mathbf z},\qquad
\rho_0=\text{constant},\qquad
p_0=\text{constant},\qquad
\mathbf u_0=0.
\end{equation}
Consider incompressible transverse perturbations varying along the field:
\begin{equation}
\mathbf u_1=u_x(z,t)\hat{\mathbf x},\qquad
\mathbf B_1=b_x(z,t)\hat{\mathbf x}.
\end{equation}
Then $\nabla\cdot\mathbf u_1=0$, $\nabla\cdot\mathbf B_1=0$, and no density or pressure perturbation is required. The displacement is perpendicular to both $\mathbf B_0$ and the parallel wave vector.


\section{Linearized induction}

The ideal induction equation is
\begin{equation}
\partial_t\mathbf B=\nabla\times(\mathbf u\times\mathbf B).
\end{equation}
To first order,
\begin{equation}
\partial_t\mathbf B_1=\nabla\times(\mathbf u_1\times\mathbf B_0).
\end{equation}
For the chosen geometry,
\begin{equation}
\boxed{\frac{\partial b_x}{\partial t}
=B_0\frac{\partial u_x}{\partial z}}.
\end{equation}
A transverse velocity gradient bends the equilibrium field and creates a transverse magnetic perturbation.


\section{Linearized momentum and magnetic tension}

The transverse momentum equation is
\begin{equation}
\rho_0\frac{\partial\mathbf u_1}{\partial t}
=\frac{1}{\mu_0}(\nabla\times\mathbf B_1)\times\mathbf B_0.
\end{equation}
For $\mathbf B_1=b_x(z,t)\hat{\mathbf x}$,
\begin{equation}
\boxed{\rho_0\frac{\partial u_x}{\partial t}
=\frac{B_0}{\mu_0}\frac{\partial b_x}{\partial z}}.
\end{equation}
This is magnetic tension: curvature of the perturbed field produces a restoring force along the transverse displacement.


\section{Wave equation and Alfvén speed}

Differentiate the momentum equation in time and use induction:
\begin{align}
\rho_0\frac{\partial^2u_x}{\partial t^2}
&=\frac{B_0}{\mu_0}\frac{\partial}{\partial z}
\left(\frac{\partial b_x}{\partial t}\right)\\
&=\frac{B_0^2}{\mu_0}\frac{\partial^2u_x}{\partial z^2}.
\end{align}
Therefore
\begin{equation}
\boxed{\frac{\partial^2u_x}{\partial t^2}
=v_A^2\frac{\partial^2u_x}{\partial z^2}},
\qquad
\boxed{v_A=\frac{B_0}{\sqrt{\mu_0\rho_0}}}.
\end{equation}
A plane wave $u_x=\tilde u e^{i(k_\parallel z-\omega t)}$ gives
\begin{equation}
\boxed{\omega^2=k_\parallel^2v_A^2}.
\end{equation}


\section{Amplitude and phase relation}

For a plane wave, induction gives
\begin{equation}
-i\omega\tilde b=ik_\parallel B_0\tilde u.
\end{equation}
Thus
\begin{equation}
\tilde b=-\frac{k_\parallel B_0}{\omega}\tilde u
=\mp\sqrt{\mu_0\rho_0}\,\tilde u,
\end{equation}
where the sign depends on propagation direction and phase convention. The velocity and magnetic perturbations are correlated; changing the propagation direction changes their relative sign.


\section{Energy density and flux}

The quadratic kinetic and magnetic perturbation energies are
\begin{equation}
\mathcal E_K=\frac{\rho_0u_1^2}{2},\qquad
\mathcal E_B=\frac{b_1^2}{2\mu_0}.
\end{equation}
Using $|b_1|=\sqrt{\mu_0\rho_0}|u_1|$ for a traveling wave,
\begin{equation}
\boxed{\mathcal E_K=\mathcal E_B}.
\end{equation}
Multiplying the momentum and induction equations by $u_x$ and $b_x/\mu_0$ yields
\begin{equation}
\frac{\partial}{\partial t}\left(
\frac{\rho_0u_x^2}{2}+\frac{b_x^2}{2\mu_0}\right)
+\frac{\partial}{\partial z}\left(-\frac{B_0u_xb_x}{\mu_0}\right)=0.
\end{equation}
The sign of the flux tracks propagation direction.


\section{Elsasser variables}

Define
\begin{equation}
\mathbf z^\pm=\mathbf u_1\mp\frac{\mathbf B_1}{\sqrt{\mu_0\rho_0}}.
\end{equation}
The linear equations become
\begin{equation}
\partial_t\mathbf z^\pm\pm v_A\partial_z\mathbf z^\pm=0.
\end{equation}
$\mathbf z^+$ and $\mathbf z^-$ describe disturbances propagating in opposite directions along the equilibrium field. This form is central in Alfvénic turbulence.


\section{Limits of the elementary derivation}

The result assumes uniform equilibrium, ideal conductivity, scalar pressure, no finite-Larmor-radius effects, and propagation exactly along $\mathbf B_0$. Nonuniform density creates an Alfvén continuum; toroidal geometry couples poloidal harmonics; kinetic particles add damping or drive. Those extensions motivate Basics 34-37 and Modules 4-6.


\section{Computational laboratory}
Integrate the coupled linear induction and momentum equations for a periodic domain. Initialize a single traveling wave and a standing wave. Verify the measured speed, amplitude relation, energy equipartition, and convergence of total perturbation energy in both Python and MATLAB.

\section{Summary}
\begin{itemize}[leftmargin=1.6em]
\item Shear-Alfvén waves are transverse ideal-MHD disturbances restored by magnetic tension.
\item Their speed is $v_A=B_0/\sqrt{\mu_0\rho_0}$ and propagation is along the field.
\item Velocity and magnetic perturbations have a direction-dependent phase relation.
\item Traveling waves have equal kinetic and magnetic perturbation energies.
\item Nonuniformity, geometry, and kinetic particles extend this elementary wave into continua and eigenmodes.
\end{itemize}

\SymbolsAcronymsSection
\begin{symtable}
$v_A$ & Alfvén speed & $B_0/\sqrt{\mu_0\rho_0}$. \\
$k_\parallel$ & parallel wave number & Phase variation along the equilibrium field. \\
$u_x$ & transverse velocity & Shear-flow perturbation. \\
$b_x$ & transverse magnetic perturbation & Field-line bending amplitude. \\
$\mathbf z^\pm$ & Elsasser variables & Counterpropagating Alfvénic combinations. \\
$\mathcal E_K,\mathcal E_B$ & perturbation energies & Kinetic and magnetic contributions. \\
\end{symtable}

\section{Exercises}
\begin{enumerate}[leftmargin=1.8em]
\item Derive the linear induction equation for the chosen polarization.
\item Derive the transverse Lorentz-force term from $(\nabla\times\mathbf B_1)\times\mathbf B_0$.
\item Obtain the Alfvén wave equation and check dimensions of $v_A$.
\item Derive the velocity-magnetic amplitude relation for both propagation directions.
\item Derive the quadratic energy conservation law.
\end{enumerate}

\section{References}
\begin{thebibliography}{99}
\bibitem{ref1} F. F. Chen, \emph{Introduction to Plasma Physics and Controlled Fusion}, 3rd ed., Springer (2016).
\bibitem{ref2} R. J. Goldston and P. H. Rutherford, \emph{Introduction to Plasma Physics}, CRC Press (1995).
\bibitem{ref3} P. M. Bellan, \emph{Fundamentals of Plasma Physics}, Cambridge University Press (2006).
\bibitem{ref4} J. P. Freidberg, \emph{Ideal MHD}, Cambridge University Press (2014).
\bibitem{ref5} J. P. Goedbloed, R. Keppens, and S. Poedts, \emph{Advanced Magnetohydrodynamics}, Cambridge University Press (2010).
\bibitem{ref6} H. Goedbloed and S. Poedts, \emph{Principles of Magnetohydrodynamics}, Cambridge University Press (2004).
\bibitem{ref7} T. H. Stix, \emph{Waves in Plasmas}, American Institute of Physics (1992).
\bibitem{ref8} D. A. Gurnett and A. Bhattacharjee, \emph{Introduction to Plasma Physics}, 2nd ed., Cambridge University Press (2017).
\bibitem{ref9} G. B. Arfken, H. J. Weber, and F. E. Harris, \emph{Mathematical Methods for Physicists}, 7th ed., Academic Press (2013).
\bibitem{ref10} G. Strang, \emph{Linear Algebra and Its Applications}, 4th ed., Brooks/Cole (2006).
\bibitem{ref11} W. A. Strauss, \emph{Partial Differential Equations: An Introduction}, 2nd ed., Wiley (2007).
\end{thebibliography}

\end{document}
